% =========================================================
\section{Giai Đoạn Mô Hình Hóa và Thực Thi}
% =========================================================

Giai đoạn cuối hiện thực hóa kiến trúc mạng nơ-ron đồ thị thời gian (Mục~\ref{sec:encoders}), lựa chọn siêu tham số, và điều phối huấn luyện. Giai đoạn gồm hai tác tử: $\mathcal{A}_{gnn}$ sinh kịch bản huấn luyện, $\mathcal{A}_{ops}$ thực thi và gỡ lỗi kịch bản đó.

\subsection{Tác Tử Chuyên Gia GNN ($\mathcal{A}_{gnn}$)}

$\mathcal{A}_{gnn}$ nhận định nghĩa bài toán từ $\mathcal{S}_{spec}$ cùng cấu trúc đồ thị $\mathcal{G}$, từ đó lựa chọn siêu tham số và sinh kịch bản huấn luyện hoàn chỉnh. Khác với các phương pháp tối ưu siêu tham số dựa trên huấn luyện thử nghiệm lặp (Bayesian Optimization~\cite{feurer2015autosklearn}, ASHA~\cite{liaw2018tune}), $\mathcal{A}_{gnn}$ áp dụng cơ chế \textit{tối ưu siêu tham số không cần huấn luyện} (training-free HPO) thông qua giao thức Model Context Protocol (MCP).

\paragraph{Tối ưu siêu tham số qua tổng hợp tri thức.}
$\mathcal{A}_{gnn}$ truy vấn năm nguồn tri thức dị thể qua các máy chủ MCP: (i)~quy tắc suy luận dựa trên đặc tính bài toán (số nút, số bảng, loại tác vụ), (ii)~bài báo khoa học từ Google Scholar và arXiv, (iii)~nghiên cứu có trích dẫn cao từ Semantic Scholar, và (iv)~giải pháp cuộc thi và notebook trên Kaggle. Các cấu hình thu thập được tổng hợp theo chiến lược \textit{bỏ phiếu trung vị} (ensemble median voting): với mỗi siêu tham số, hệ thống lấy giá trị trung vị từ tập đề xuất, giảm ảnh hưởng của ngoại lai từ nguồn đơn lẻ.

Quá trình truy vấn MCP được bảo vệ bởi cơ chế ngắt mạch (circuit breaker): sau ba lần thất bại liên tiếp, lời gọi đến nguồn đó bị tạm dừng. Khi tất cả nguồn không khả dụng, hệ thống quay về bộ siêu tham số mặc định.

\paragraph{Lựa chọn độ đo đánh giá.}
$\mathcal{A}_{gnn}$ ưu tiên độ đo do người dùng chỉ định; nếu không có, áp dụng bảng ánh xạ mặc định theo loại tác vụ (MAE cho hồi quy, average precision cho phân loại nhị phân, accuracy cho phân loại đa lớp, MAP cho dự đoán liên kết). Các tên gọi phổ biến (AUROC, AUC, F1) được chuẩn hóa tự động.

\paragraph{Sinh kịch bản huấn luyện từ khuôn mẫu cố định.}
Sau khi xác định siêu tham số, $\mathcal{A}_{gnn}$ sinh kịch bản huấn luyện từ một \textit{khuôn mẫu mã cố định} (code template) tuân thủ quy trình tham chiếu của RelBench~\cite{robinson2024relbench}. Kịch bản hiện thực hóa kiến trúc bốn thành phần đã mô tả trong Mục~\ref{sec:encoders}: mã hóa đặc trưng nút dị thể, mã hóa thời gian tương đối, truyền thông điệp GraphSAGE~\cite{hamilton2017inductive} trên đồ thị dị thể, và đầu dự đoán MLP. Kịch bản cũng bao gồm cấu hình hàm mất mát theo loại tác vụ, bộ nạp dữ liệu với lấy mẫu lân cận nhất quán thời gian (Thuật toán~\ref{alg:computation_graph}), vòng lặp huấn luyện với dừng sớm, và ghi kết quả đánh giá.

Thiết kế khuôn mẫu cố định --- thay vì để LLM tự sinh toàn bộ mã huấn luyện --- loại bỏ rủi ro LLM sinh mã có lỗi logic tinh vi (rò rỉ dữ liệu, sai thứ tự đánh giá, thiếu lưu mô hình). System prompt cấm tác tử viết kịch bản thủ công và bắt buộc sử dụng khuôn mẫu.

\subsection{Tác Tử Vận Hành ($\mathcal{A}_{ops}$)}
\label{sec:operation-agent}

$\mathcal{A}_{ops}$ thực thi kịch bản huấn luyện do $\mathcal{A}_{gnn}$ sinh ra. Tác tử này hoạt động theo thiết kế \textit{lai} (hybrid), lấy ý tưởng từ Self-Debugging~\cite{chen2024selfdebugging} và AutoML-Agent~\cite{trirat2025automl}: đường thực thi chính không sử dụng LLM; LLM chỉ được kích hoạt khi kịch bản thất bại.

\paragraph{Pha 1: Thực thi tất định.}
$\mathcal{A}_{ops}$ khởi tạo tiến trình con cô lập để chạy kịch bản huấn luyện, thu thập cả đầu ra chuẩn và đầu ra lỗi theo thời gian thực. Nếu kịch bản hoàn tất mà không có lỗi, tác tử chuyển trực tiếp sang trích xuất kết quả --- không tiêu tốn token LLM nào trên đường chính.

\paragraph{Pha 2: Gỡ lỗi tự động.}
Khi kịch bản thất bại, $\mathcal{A}_{ops}$ kích hoạt vòng lặp gỡ lỗi tối đa $N$ lần ($N=3$). Mỗi vòng gồm bốn bước: (1)~khôi phục toàn bộ tệp mã về bản gốc, ngăn tích lũy lỗi từ bản vá thất bại trước; (2)~phân tích traceback để xác định tệp nguồn lỗi --- có thể là tệp định nghĩa đồ thị hoặc tệp tác vụ, không nhất thiết là kịch bản huấn luyện; (3)~gửi mã lỗi, traceback, và lịch sử các lần gỡ lỗi trước cho LLM để sinh bản vá; (4)~thực thi lại kịch bản.

Nếu cả $N$ vòng thất bại, $\mathcal{A}_{ops}$ khôi phục tệp gốc trước khi trả quyền cho tầng điều phối (Mục~\ref{sec:self-recovery}). Việc khôi phục là điều kiện cần để tầng điều phối có thể sinh lại mã từ đầu mà không bị ảnh hưởng bởi bản vá hỏng.

\paragraph{Sản phẩm đầu ra.}
Khi huấn luyện hoàn tất, pipeline sinh ra năm sản phẩm: định nghĩa đồ thị, định nghĩa tác vụ, kịch bản huấn luyện, trọng số mô hình đã huấn luyện, và kết quả đánh giá. $\mathcal{A}_{ops}$ cũng sinh mã suy luận (inference code) để người dùng thực hiện dự đoán trên dữ liệu mới.
